
\section{Related Work}

\subsection{Self-evolving agent}
The research focus in artificial intelligence is undergoing a profound paradigm shift: from developing static foundational models to building dynamic, self-evolving agents capable of continuous adaptation and learning~\citep{gao2025survey}. To achieve this goal, researchers are exploring various approaches. For example, some works~\citep{zhou2022large, wang2023promptagent, pryzant2023automatic} abstract the prompt generation process into a black-box optimization problem, systematically searching for and optimizing instructions to maximize the performance of large language models (LLMs) on specific tasks. Regarding agent capability building, some cutting-edge work draws on concepts from cognitive science, helping agents accumulate skills and experience through course learning or free exploration, forming a reusable skill base~\citep{wang2023voyager, wu2024copilot, qian2024iterative} or optimized toolsets~\citep{tang2025agent, qiu2025alita}. Another important technical approach is to empower agents with the ability to self-reflect and iterate. By introducing language feedback mechanisms and comparing and reflecting with ground truth answers, agents can continuously review and strengthen their decision-making logic and action capabilities~\citep{shinn2023reflexion, liang2025sage}.

\subsection{LLM Agent Memory Mechanisms}
Research on memory mechanisms for LLM agents aims to enable them to store, retain, and recall past experiences, facilitating the transition from simple reactive models to advanced agents capable of maintaining context and autonomous adaptation. Research in this domain often draws inspiration from human cognitive models, classifying memory into short-term working memory for immediate task processing and long-term memory for persistent learning~\citep{krishnan2025ai}, which relies on external storage such as vector databases~\citep{douze2024faiss} and knowledge graphs~\citep{webber2012programmatic}. To address the challenges of information retrieval and potential information flooding arising from accumulating long-term memories, one line of research focuses on optimizing memory management and structure. For example, Mem0~\citep{chhikara2025mem0} implements precise control over memory content by defining explicit memory operations, while MemInsight~\citep{salama2025meminsight} enhances semantic information by augmenting raw memories with summaries and tags to optimize subsequent retrieval efficiency. Another branch of research focuses on constructing procedural memory by generalizing reusable experiences and workflows from agents' historical execution trajectories. Specifically, ExpeL~\citep{zhao2024expel} collects execution trajectories and refines them into natural language insights and rules. Agent Workflow Memory~\citep{wang2024agent} focuses on generalizing reusable workflows from individual experiences. Memp~\citep{fang2025memp} aims to build a learnable, updatable, and lifelong procedural memory, allowing agents to acquire skills and habits through experience. Although these advanced memory mechanisms are validated on various text-based benchmarks~\citep{yang2018hotpotqa, shridhar2020alfworld, yao2022webshop, thorne2018fever} and web agent benchmarks~\citep{zhou2023webarena, deng2023mind2web}, existing test environments often lack sufficient complexity and long-term dependency requirements. Consequently, they may not fully assess the true efficacy of these mechanisms in handling complex, long-horizon, real-world tasks.